\documentclass[11pt,a4paper]{article}
\usepackage[T1]{fontenc}
\usepackage[utf8]{inputenc}
\usepackage{lmodern}
\usepackage[margin=24mm]{geometry}
\usepackage{amsmath,amssymb,bm,mathtools}
\usepackage{booktabs,longtable,array}
\usepackage{microtype}
\usepackage{xcolor}
\usepackage[colorlinks=true,linkcolor=blue!45!black,citecolor=blue!45!black,urlcolor=blue!45!black]{hyperref}
\usepackage{enumitem}
\setlist{itemsep=3pt,topsep=5pt}
\setlength{\parskip}{4pt}
\setlength{\parindent}{0pt}
\allowdisplaybreaks[2]
\numberwithin{equation}{section}
\newcommand{\dd}{\,\mathrm d}
\newcommand{\dev}{\operatorname{dev}}
\newcommand{\sym}{\operatorname{sym}}
\newcommand{\tr}{\operatorname{tr}}
\newcommand{\sgn}{\operatorname{sgn}}
\newcommand{\Div}{\operatorname{div}}
\newcommand{\E}{\bm E}
\newcommand{\Ev}{\bm E^{v}}
\newcommand{\D}{\bm D}
\newcommand{\Sg}{\bm S}
\newcommand{\T}{\bm T}
\newcommand{\ta}{\bm\tau}
\newcommand{\eps}{\bm\varepsilon}
\newcommand{\R}{\mathcal R}
\newcommand{\calE}{\mathcal E}
\newcommand{\F}{\mathcal F}
\newcommand{\ind}{\mathcal I}
\newcommand{\norm}[1]{\left\lVert#1\right\rVert}
\newcommand{\avg}[1]{\left\langle#1\right\rangle}
\title{\textbf{Bulk cohesive shear damage and its surface reduction}\\[5pt]
\large Constitutive potential, dissipation, intact activation,\\
and the connection to the reconstructed-fault model}
\author{Technical derivation for the phase-field fault project}
\date{20 October 2026}

\begin{document}
\maketitle
\begin{abstract}
The finite-tip obstruction in P4a is caused by restricting the mechanical
damage mechanism to normal profiles attached to an already reconstructed
surface. This note constructs a bulk-defined, small-strain constitutive model
with recoverable cohesive shear, bulk and cohesive Maxwell relaxation,
constant Coulomb friction, and irreversible phase evolution. Its energy
balance and intact limit are derived explicitly. The existing surface storage
$Z^2/(2GI)$ is recovered by a normal-profile reduction. An exact residual
decomposition identifies the energy, dissipation, and history lost if that
reduction is imposed where it is not valid. In particular, equality of
instantaneous energies does not establish equivalence of evolving-profile
constitutive laws. A bulk-authoritative connection is proposed that permits
damage ahead of the surface without adding duplicate cohesive energy.
\end{abstract}

\textbf{Status.} This is a constitutive proposal and analytical derivation,
not a claim that the production implementation already solves these equations.
It completes the requested derivation for a fixed-direction P4 setting.
The new bulk model is not uniquely implied by the old surface equations.
Its explicitly selected assumptions are listed below.

\tableofcontents
\newpage

\section{Scope, conventions, and the specific defect}

We select infinitesimal strain, isothermal incompressible mechanics, a fixed
material domain $\Omega$, and one prescribed shear direction. There is no
advection, spin, moving material mesh, opening, branching, or RSF state in
this derivation. Time derivatives are material derivatives in this fixed
small-strain description. The surface may be reconstructed later as a
description of the bulk state; no physical motion or creation of that surface
is required to evaluate the bulk constitutive law.

Let $\bm u$ denote displacement, $\bm v=\dot{\bm u}$ velocity, and
\begin{equation}
 \E=\dev\sym\nabla\bm u,\qquad
 \D=\dot\E=\dev\sym\nabla\bm v,\qquad
 \Div\bm v=0.
\end{equation}
The constant unit vectors $\bm m,\bm n$ are perpendicular, and
\begin{equation}
 \Sg=\tfrac12(\bm m\otimes\bm n+\bm n\otimes\bm m),\qquad
 \Sg:\Sg=\tfrac12,\qquad \tr\Sg=0.
\end{equation}
For the horizontal example $S_{xy}=S_{yx}=1/2$. The engineering shear is
$\gamma=2\E:\Sg$ and the resolved shear stress is $t=\ta:\Sg$.
All tensor contractions are three-dimensional. For a trace-free plane-strain
stress,
\begin{equation}
 \norm{\ta}_3^2=\tau_{xx}^2+\tau_{yy}^2+
                 (\tau_{xx}+\tau_{yy})^2+2\tau_{xy}^2.
\end{equation}
Pressure $p$ is a Lagrange multiplier: $\bm\sigma=-p\bm 1+\ta$.
The prescribed frictional normal compression $p_N\geq0$ is distinct from $p$.
We take constant residual resistance $\tau_r=\mu p_N\geq0$.

For $0\leq\phi<1$, choose
\begin{equation}
 h(\phi)=g(\phi)^{-1}-1,\qquad h(0)=0,\qquad h'\geq0,
 \qquad
 h(\phi)=\frac{m_g\phi(1+p_g\phi)}{(1-\phi)^2}
 \label{eq:h}
\end{equation}
for the retained Lorentz-type choice. Subscripts distinguish the parameters
$m_g,p_g$ from direction and pressure. The fracture functional is
\begin{equation}
 \F[\phi]=K_f\int_\Omega
       \bigl(\phi+\ell^2|\nabla\phi|^2\bigr)\dd x,\qquad
 K_f=\frac{G_c}{c_0\ell},\qquad c_0=\frac83.
 \label{eq:fracture}
\end{equation}
It is counted once. The fully damaged endpoint $g=0$ requires its own limiting
treatment and is not regularized by an undocumented floor.

P4a uses a strip indicator in
$\ta=\T-2Z\mathbf1_{A_\Gamma}(h/I)\Sg$.
A phase perturbation supported wholly outside $A_\Gamma$ changes neither
this stress nor the existing surface integral $I$.
Thus $\delta(E_b+E_c)=0$ there, regardless of load magnitude. This is a
support restriction of the constitutive model, not an integration error.
The model below removes that restriction by defining its mechanical variables
throughout $\Omega$, including where $\phi=0$ and where $\Gamma$ does not exist.

\section{Bulk state and free energy}

\subsection{State variables and kinematics}

Introduce the bulk viscous deviatoric strain $\Ev$, the signed additional
shear deformation $k$, and its viscous part $b$. All are bulk fields.
The recoverable cohesive shear is
\begin{equation}
 a=k-b,\qquad
 \bm e=\E-\Ev-k\Sg.
 \label{eq:kinematics}
\end{equation}
Here $a,b,k$ are dimensionless engineering shear measures, not slips in metres.
The local slip-deformation rate is $w=\dot k$, with units $\mathrm{s}^{-1}$.
The signed surface slip rate introduced later has units $\mathrm{m\,s}^{-1}$.

The stored mechanical energy density is selected as
\begin{equation}
 \boxed{\psi(\E,\Ev,k,b,\phi)
       =G\,\bm e:\bm e+\frac{G(k-b)^2}{2h(\phi)}.}
 \label{eq:bulkenergy}
\end{equation}
The first term is bulk elastic storage; the second is recoverable cohesive
storage. There is no additional surface storage in the bulk formulation.
The total energy is
$\calE=\int_\Omega\psi\dd x+\F[\phi]$.
For now $G>0$ is constant and the same modulus appears in both terms.

The zero-compliance limit is an extended-valued energy:
\begin{equation}
 \frac{Ga^2}{2h}:=
 \begin{cases}
  Ga^2/(2h),&h>0,\\
  0,&h=0,\ a=0,\\
  +\infty,&h=0,\ a\ne0.
 \end{cases}
 \label{eq:closedperspective}
\end{equation}
Consequently intact material admits no finite cohesive elastic deformation.
This definition does not require a fictitious initial damage or a positive
normalization floor.

\subsection{Conjugate forces}

For $h>0$, differentiation gives
\begin{align}
 \ta&=2G\bm e,&
 q&=\frac{Ga}{h},&
 t&=\ta:\Sg,&
 r&=t-q, \label{eq:forces}\\
 \partial_{\Ev}\psi&=-\ta,&
 \partial_k\psi&=-t+q,&
 \partial_b\psi&=-q,&
 \partial_\phi\psi&=-\frac{h'q^2}{2G}.
 \label{eq:partials}
\end{align}
The symbol $q$ is a \emph{bulk} cohesive traction. It is not defined by
division by a surface normalization. A surface traction $q_\Gamma$ will be
introduced by reduction.

For fixed incoming mechanical state, the mechanical phase driving is
\begin{equation}
 Y_{\rm mech}=-\partial_\phi\psi
       =\frac{h'q^2}{2G}
       =-g'\frac{q^2}{2Gg^2}.
 \label{eq:driving}
\end{equation}
This identity does not introduce a maximum-history variable $H$.
Irreversibility is imposed directly on the physical phase field.
The intact value of \eqref{eq:driving} is a one-sided relaxed limit,
derived in Section~\ref{sec:intact}; arbitrary evaluation of $0/0$ is excluded.

\section{Dissipation potential and constitutive laws}

\subsection{A selected generalized-standard-material structure}

Choose equal bulk and cohesive relaxation times by using the same
viscosity $\eta>0$, and put $\lambda=\eta/G$. With state held fixed, define
the dissipation-potential density
\begin{equation}
 \boxed{\R=
  \eta\,\dot\Ev:\dot\Ev
  +\frac{\eta}{2h}\dot b^2
  +\tau_r|w|
  +\frac{\beta_\phi}{2}\dot\phi^2
  +\ind_{[0,\infty)}(\dot\phi),\qquad \beta_\phi\geq0.}
 \label{eq:rayleigh}
\end{equation}
The viscous perspective at $h=0$ is zero for $\dot b=0$ and infinite otherwise.
The indicator is zero on its stated set and infinite outside it.
For the retained rate-independent phase model, set $\beta_\phi=0$.
It is displayed to identify optional phase viscosity, not to recommend adding it.

Normality with respect to the \emph{rates} gives
\begin{align}
 \dot\Ev&=\frac{\ta}{2\eta},&
 \dot b&=\frac{h q}{\eta}=\frac a\lambda,&
 r&\in\tau_r\partial|w|.
 \label{eq:flow}
\end{align}
Thus
\begin{equation}
 |r|\leq\tau_r \quad (w=0),\qquad
 r=\tau_r\sgn w \quad (w\ne0),\qquad rw=\tau_r|w|.
 \label{eq:coulomb}
\end{equation}
This is a two-sided Coulomb inclusion. Restricting to positive slip gives the
usual complementarity $w\geq0$, $\tau_r-r\geq0$,
$w(\tau_r-r)=0$.
No accumulated frictional work is added to stored energy.

The state dependence of $\R$ matters: constitutive normality differentiates
$\R$ with respect to rates while holding the current state fixed.
Differentiating $\R$ with respect to $\phi$ as though it were free energy
would define a different model. A time-discrete minimization with a
phase-dependent viscous metric must be derived consistently; this note does
not assert equivalence to the existing endpoint algorithm at finite time step.

\subsection{Phase inequality and the global power balance}

For homogeneous natural phase boundary conditions,
\begin{equation}
 \mathcal Y=\frac{h'q^2}{2G}-K_f(1-2\ell^2\Delta\phi).
 \label{eq:Ytotal}
\end{equation}
Away from the upper endpoint, phase normality is
\begin{equation}
 \dot\phi\geq0,\qquad
 \beta_\phi\dot\phi-\mathcal Y\geq0,\qquad
 \dot\phi(\beta_\phi\dot\phi-\mathcal Y)=0.
 \label{eq:phaseVI}
\end{equation}
Equivalently, use a weak variational inequality and the accepted lower obstacle
$\phi_{n+1}\geq\phi_n$ in this nonadvecting material configuration.
The gradient term can make the driving nonlocal even though the constitutive
traction is defined at every bulk point.

Differentiating \eqref{eq:bulkenergy}, without first imposing flow laws,
\begin{equation}
 \dot\psi=\ta:\D-\ta:\dot\Ev-(t-q)w-q\dot b
                   -\frac{h'q^2}{2G}\dot\phi.
 \label{eq:localpower}
\end{equation}
For quasistatic equilibrium
$\Div(\ta-p\bm1)+\bm f=0$ and incompressibility,
$\int_\Omega\ta:\D\dd x=P_{\rm ext}$, with prescribed-boundary reactions
included in $P_{\rm ext}$. Combining \eqref{eq:flow}--\eqref{eq:phaseVI}
and integrating the phase-gradient term by parts gives
\begin{equation}
 \boxed{\dot{\calE}=P_{\rm ext}
 -\int_\Omega\left[
   \frac{\norm{\ta}_3^2}{2\eta}
   +\frac{h q^2}{\eta}
   +\tau_r|w|
   +\beta_\phi\dot\phi^2
 \right]\dd x.}
 \label{eq:balance}
\end{equation}
Every displayed dissipative term is nonnegative. Nonzero prescribed phase
flux contributes the corresponding boundary microtraction power.
Equation~\eqref{eq:balance} is a smooth-path identity. Nonconvex
rate-independent jumps require an energetic solution/selection and an energy
inequality; stationarity alone does not prove global stability or determine
unresolved jump dissipation.

\section{Independent bulk history and compatibility with Maxwell mechanics}

Introduce
\begin{equation}
 \bm A=\Ev+b\Sg,\qquad
 \T=2G(\E-\bm A),\qquad
 \boxed{\ta=\T-2Ga\Sg.}
 \label{eq:T}
\end{equation}
For constant $G,\eta,\Sg$, \eqref{eq:flow} implies
\begin{equation}
 \dot{\bm A}=\frac{\ta}{2\eta}+\frac a\lambda\Sg
             =\frac{\T}{2\eta},
 \qquad
 \boxed{\dot\T+\lambda^{-1}\T=2G\D,\quad
        \dot a+\lambda^{-1}a=w.}
 \label{eq:Ta}
\end{equation}
Hence the actual bulk stress satisfies
\begin{equation}
 \boxed{\dot\ta+\lambda^{-1}\ta=2G(\D-w\Sg).}
 \label{eq:maxwellbulk}
\end{equation}
The independent history $\T$ has the same simple Maxwell form used in P2.
The additional mechanical variable is now the \emph{bulk} cohesive elastic
shear $a$, not only a surface moment.

The mechanical storage may equivalently be evaluated as
\begin{equation}
 \psi(\T,a,\phi)=
       \frac{\norm{\T-2Ga\Sg}_3^2}{4G}
       +\frac{Ga^2}{2h}.
 \label{eq:Taenergy}
\end{equation}
At fixed $\T,a$, only the second term has explicit phase dependence.
This agrees with varying \eqref{eq:bulkenergy} at fixed
$\E,\Ev,k,b$. It differs from fixing a surface moment while changing its
bulk profile, as shown in Section~\ref{sec:evolving}.

The convenient cancellation in \eqref{eq:Ta} uses equal relaxation times.
For distinct $\lambda_b=\eta_b/G$, $\lambda_c=\eta_c/G$,
\begin{equation}
 \dot\T+\lambda_b^{-1}\T
  =2G\D+2Ga(\lambda_b^{-1}-\lambda_c^{-1})\Sg.
\end{equation}
Neither variable coefficients nor objective transport can be appended without
their product-rule terms.

\section{The intact limit and a definite activation criterion}
\label{sec:intact}

\subsection{Mixed traction form}

The cohesive perspective admits the exact dual representation
\begin{equation}
 \frac{Ga^2}{2h}
   =\sup_q\left(qa-\frac{h q^2}{2G}\right),\qquad h\geq0,
 \label{eq:dual}
\end{equation}
including the extended values in \eqref{eq:closedperspective}.
Thus a mixed density is
\begin{equation}
 \psi_{\rm mix}=G\bm e:\bm e+q(k-b)-\frac{h q^2}{2G}.
 \label{eq:mixed}
\end{equation}
The saddle variable $q$ enforces $k-b=hq/G$.
At $h=0$, $k=b$, while $q$ can be finite. It acts as a constraint traction,
not as a quotient of two prescribed zeros.
The energy is the supremum value, not the mixed expression evaluated at an
arbitrary $q$.

With friction, intact sticking alone does not uniquely select $q$:
$|t-q|\leq\tau_r$ is an interval.
Consequently the saddle equations alone are not a complete numerical birth
rule. We select the one-sided relaxed onset below, which fixes the relevant
virgin limit from the same energy and frictional dissipation.

\subsection{A virgin elastic shear calculation}

Freeze the viscous histories during an infinitesimal damage event, take one
shear mode, and let $\gamma$ be its available elastic trial shear.
Write $t_0=G\gamma$ for its intact trial stress. Starting with zero cohesive
history, the work required to reach $a$ is
\begin{equation}
 J(a,\phi;\gamma)
   =\frac G2(\gamma-a)^2+\frac{Ga^2}{2h(\phi)}
       +\tau_r|a|.
 \label{eq:Jvirgin}
\end{equation}
The last term is dissipated frictional work for a monotonic virgin change;
it is not stored energy and is not a formula for arbitrary cyclic loading.
Minimization gives, with
$z=\sgn(t_0)(|t_0|-\tau_r)_+$,
\begin{align}
 a_*&=\frac{h}{G(1+h)}z
       =\frac{1-g}{G}z,&
 q_*&=g z, \label{eq:soft}\\
 J_{\rm red}(\gamma,\phi)
     &=\frac12G\gamma^2-\frac{1-g(\phi)}{2G}z^2,&
 \partial_\phi J_{\rm red}&=\frac{g'}{2G}z^2.
 \label{eq:Jred}
\end{align}
For positive sliding this yields
$t=\tau_r+g(t_0-\tau_r)$.
For $|t_0|\leq\tau_r$, $a_*=q_*=0$ and the response sticks.
In the virgin limit,
\begin{equation}
 a_*\to0,\qquad q_*\to z,\qquad
 -\partial_{\phi+}J_{\rm red}(0)=\frac{h'(0)}{2G}z^2.
 \label{eq:intactforce}
\end{equation}
Finite traction with zero cohesive deformation is therefore consistent.
The selected multiplier is the minimum-magnitude member of the intact
friction interval, obtained by the $h\downarrow0$ relaxed limit.
Holding $a=0$ for every positive trial $\phi$ instead would remove this
relaxation and give a zero derivative.

For homogeneous virgin material, the AT1 onset condition is
\begin{equation}
 \frac{m_g}{2G}(|t_0|-\tau_r)_+^2>K_f.
 \label{eq:onset}
\end{equation}
Equivalently,
\begin{equation}
 H_t=\frac{K_f}{m_g},\qquad
 \boxed{\tau_p=\tau_r+\sqrt{2G H_t}.}
 \label{eq:strength}
\end{equation}
The equality is onset neutrality; below it the lower obstacle is stable.
This is a local homogeneous calibration, not the solution of a spatial
crack-tip boundary-value problem.

For the P2 constants $G=10^6$ Pa, $\ell=0.15625$ m,
$G_c=26.6666666667$ J/m$^2$, $m_g=128$, and $\tau_r=1000$ Pa:
$K_f\simeq64$ Pa, $H_t\simeq0.5$ Pa and $\tau_p\simeq2000$ Pa.
The pre-damaged P2 band has a different onset problem.
Also, P4a's reported trial shear of roughly $0.806$--$0.827$ kPa in its tip
window is below this example's residual resistance: a repaired law need not
produce growth in that prescribed snapshot. The demonstrated defect is the
structural absence of driving for \emph{any} load outside the old strip.

\subsection{Relationship to established frictional phase-field laws}

The positive-slip specialization above has the stress interpolation and
peak/residual strength structure used by Fei and Choo~\cite{fei2020}.
Their construction accounts for friction in fracture driving; it is not
obtained by degrading all compressive elastic energy.
The later in-plane formulation~\cite{fei2026} explicitly evaluates strength
and activation in initially intact material.
These precedents support defining a bulk activation law independent of an
existing reconstructed surface. They do not prove that the present Maxwell
extension, its unloading law, or its normalized surface reduction is identical
to those published models.

\section{Reduction to a flat reconstructed surface}
\label{sec:reduction}

\subsection{Assumptions needed for an exact energy identity}

Take a fixed straight surface parameterized by arclength $s$ and Cartesian
normal coordinate $n$, so $\dd x=\dd s\,\dd n$.
Profiles do not overlap, $\Sg$ is constant on each profile, and $G,\eta,\tau_r$
are constant across it. The integration interval is fixed and includes the
represented normal support. Define
\begin{equation}
 I(s)=\int h(\phi(s,n))\dd n>0,\qquad
 \chi(s,n)=\frac{h(\phi(s,n))}{I(s)},\qquad
 \int\chi\dd n=1.
 \label{eq:I}
\end{equation}
The surface elastic slip and moment are
\begin{equation}
 s_e(s)=\int a(s,n)\dd n,\qquad Z(s)=G s_e(s).
 \label{eq:moment}
\end{equation}
The profile ansatz is
\begin{equation}
 a(s,n)=\chi(s,n)s_e(s).
 \label{eq:ansatz}
\end{equation}
It implies constant cohesive traction on a profile,
$q(s,n)=q_\Gamma(s)=Z(s)/I(s)$.
It does not imply that the actual shear stress $t(s,n)$ is constant.

Substitution into the \emph{bulk} energy gives
\begin{align}
 \ta&=\T-2Z\chi\Sg,\\
 \int\frac{Ga^2}{2h}\dd n
    &=\frac{G s_e^2}{2I^2}\int h\dd n
      =\boxed{\frac{G s_e^2}{2I}=\frac{Z^2}{2GI}}.
 \label{eq:surfaceenergy}
\end{align}
Thus the current cohesive surface term is a reduced representation of the
bulk cohesive storage. Adding it to that same unreduced bulk storage would
count the energy twice.

The ansatz also minimizes \emph{cohesive storage alone} for a given moment:
varying $\int Ga^2/(2h)\dd n$ subject to $\int a\dd n=s_e$ gives $Ga/h$
constant. It does not generally minimize the full bulk-plus-cohesive energy,
whose $a$ derivative also contains $-t$.

\subsection{The weak traction and fixed-profile dissipation}

At fixed $\phi,\T$, variation with respect to $s_e$ gives
\begin{equation}
 \partial_{s_e}E_{\rm mech}
     =q_\Gamma-t_\Gamma,\qquad
 t_\Gamma=\int\chi\,(\ta:\Sg)\dd n.
 \label{eq:weaktraction}
\end{equation}
This is the work-adjoint traction; it is not in general centerline sampling.
For a fixed profile and restricted rate $w=\chi V$,
\begin{equation}
 \dot s_e+\lambda^{-1}s_e=V,\qquad
 \dot Z+\lambda^{-1}Z=GV.
 \label{eq:Zlaw}
\end{equation}
The cohesive viscous dissipation and friction reduce to
\begin{align}
 \int\frac{h q^2}{\eta}\dd n
   &=\frac{Z^2}{\eta I},&
 \int\tau_r|w|\dd n&=\tau_r|V|.
 \label{eq:reducedD}
\end{align}
The corresponding restricted surface law is
\begin{equation}
 t_\Gamma-q_\Gamma\in\tau_r\partial|V|.
 \label{eq:surfflow}
\end{equation}
These equations reproduce the current straight-band energy and Maxwell
moment structure.

There are two different statements here. Restricting the fields and virtual
rates to the profile ansatz yields a consistent \emph{reduced model}.
An unrestricted bulk solution is represented exactly by it only if that
solution satisfies the ansatz and its local flow conditions. In particular,
a constant $q$ and a common slip sign require constant $t-q$ under local
Coulomb sliding. A surface weak balance alone does not impose that pointwise
condition.

\section{Evolving profiles: what changes and what does not}
\label{sec:evolving}

\subsection{The missing rate in a direct identification}

If \eqref{eq:ansatz} is maintained as phase evolves,
\begin{equation}
 \dot a=\chi\dot s_e+s_e\dot\chi.
\end{equation}
Combining this with the bulk law \eqref{eq:Ta} and the moment law
\eqref{eq:Zlaw} requires
\begin{equation}
 \boxed{w=\chi V+s_e\dot\chi.}
 \label{eq:wredistribute}
\end{equation}
Since $\int\dot\chi\dd n=0$, its normal integral remains $V$.
The additional local shear rate redistributes recoverable deformation; it is
not generally zero pointwise.
With \eqref{eq:wredistribute}, the physical stress equation becomes
\begin{equation}
 \dot\ta+\lambda^{-1}\ta
   =2G(\D-\chi V\Sg)-2G s_e\dot\chi\Sg,
 \label{eq:currentredistribution}
\end{equation}
which is the fixed-direction redistribution term of the current surface
model. Its algebraic agreement does not by itself establish agreement of
frictional dissipation.

Indeed, bulk Coulomb friction requires
\begin{equation}
 \mathcal D_{f,\mathrm{profile}}
  =\tau_r\int|\chi V+s_e\dot\chi|\dd n
  \ \geq\ \tau_r|V|.
 \label{eq:redistribfriction}
\end{equation}
Equality holds if the local rate has a common sign almost everywhere.
For $V=0$, nonzero $s_e\dot\chi$ can dissipate energy in this bulk law while
the current surface friction law dissipates zero.
This is a constitutive difference, not a quadrature defect.
The current model may instead define friction only on the collective
surface slip and treat profile redistribution as a separate recoverable
mechanism. That is a distinct constrained constitutive choice; it must not
be called an exact reduction of arbitrary local Coulomb friction.

\subsection{Why the two phase derivatives differ}

Let $\delta\phi$ change $\chi$, holding $\T,s_e$ fixed on the ansatz.
Then $\delta a=s_e\delta\chi$ and
\begin{align}
 \delta E_{\rm mech}^{\mathrm{red}}
  &= -\int_{A_\Gamma}\frac{h'q_\Gamma^2}{2G}\delta\phi\dd x
     +\int_\Gamma s_e\int
        (q_\Gamma-t)\delta\chi\dd n\dd s
       \label{eq:chain}\\
  &=-\int_\Gamma s_e\int t\,\delta\chi\dd n\dd s
    -\int_\Gamma\frac{Z^2}{2GI^2}\delta I\dd s.
       \label{eq:currentphase}
\end{align}
We used $\int\delta\chi\dd n=0$ and
$\delta I=\int h'\delta\phi\dd n$.
Equation~\eqref{eq:currentphase} is the continuum phase variation retained
by the current bulk-plus-surface energy.

The first term of \eqref{eq:chain} is the bulk constitutive phase derivative
at fixed $a$. The second is work caused by moving the ansatz as phase changes.
It vanishes for profile-constant $q_\Gamma-t$, including an isostress,
common-sign sliding specialization. It need not vanish near a tip.
This proves both the validity of the current chain rule and the absence of
a general equivalence to the unrestricted local phase derivative.

The reduced surface model is not thereby thermodynamically inconsistent:
its own phase force includes the constraint-motion work and its own chosen
surface friction can close its energy balance. Its limitation is that this
mechanism is still defined only on profiles attached to $\Gamma$.

\section{An exact bulk--surface connection with a retained residual}
\label{sec:residual}

\subsection{Moment plus unresolved normal variation}

For any finite-energy bulk field on a positive-$I$ profile, define
\begin{equation}
 \boxed{a=\chi s_e+\rho,\qquad
       s_e=\int a\dd n,\qquad \int\rho\dd n=0.}
 \label{eq:rho}
\end{equation}
Here $\rho$ is a cohesive shear residual, not a new damage variable.
Where $h=0$, finite energy requires $a=\rho=0$.
The local traction and physical stress are exactly
\begin{equation}
 q=q_\Gamma+\frac{G\rho}{h},\qquad
 q_\Gamma=\frac{G s_e}{I}
          =\frac{\int hq\dd n}{I},\qquad
 \ta=\T-2\chi Z\Sg-2G\rho\Sg.
 \label{eq:rhostress}
\end{equation}
Quotients in the first formula are used only where $h>0$; mixed limits apply
on intact material.

Expanding the cohesive energy and using $\int\rho\dd n=0$ gives the exact identity
\begin{equation}
 \boxed{\int\frac{Ga^2}{2h}\dd n
   =\frac{Z^2}{2GI}+\frac G2\int\frac{\rho^2}{h}\dd n.}
 \label{eq:variance}
\end{equation}
The cross term is $(G s_e/I)\int\rho\dd n=0$.
Dropping $\rho$ therefore discards a known nonnegative stored energy.
Preserving physical stress by changing $\T$ does not preserve that discarded
cohesive storage.

Likewise,
\begin{equation}
 \int\frac{h q^2}{\eta}\dd n
 =\frac{Z^2}{\eta I}+\frac{G^2}{\eta}\int\frac{\rho^2}{h}\dd n.
 \label{eq:rhodiss}
\end{equation}
For $V=\int w\dd n$, the frictional dissipation satisfies the exact split
\begin{equation}
 \int\tau_r|w|\dd n
 =\tau_r|V|+
   \tau_r\left(\int|w|\dd n-\left|\int w\dd n\right|\right).
 \label{eq:frictionexcess}
\end{equation}
The last term is nonnegative and measures cancellation of oppositely signed
local rates. It is already present in the bulk law; it is not an added
numerical damping term.

\subsection{Residual evolution and phase-coordinate changes}

For a fixed normal-coordinate geometry and homogeneous $\lambda$,
integration of \eqref{eq:Ta} gives
$\dot s_e+\lambda^{-1}s_e=V$ exactly, with no profile ansatz.
Then
\begin{equation}
 \boxed{\dot\rho+\lambda^{-1}\rho
       =w-\chi V-s_e\dot\chi.}
 \label{eq:rhoevolution}
\end{equation}
Thus $\rho=0$ is preserved only when the right-hand side vanishes.
With $w=\chi V$ and an evolving profile, a nonzero
$-s_e\dot\chi$ generates residual history. This is the precise obstruction to
using only a surface moment as a universally closed bulk history.

Changing variables does not change the physical phase force if derivatives
are transformed correctly. At fixed physical $a$ and fixed profile geometry,
$s_e$ is fixed but
\begin{equation}
 \delta\rho=-s_e\delta\chi.
 \label{eq:coordinatephase}
\end{equation}
Holding $\rho=0$ while varying $\phi$ instead moves the physical field $a$.
The two derivatives answer different constitutive questions.
This remains important even at an instant when the residual itself is zero.

\subsection{A connection that counts each energy once}

Let $A_\Gamma$ be a union of nonoverlapping flat-profile patches.
Keep the bulk variables authoritative throughout $\Omega$.
On $A_\Gamma$, replace cohesive bulk storage by its exact representation
\eqref{eq:variance}; outside it retain the original bulk storage:
\begin{align}
 \calE={}&\int_\Omega\frac{\norm{\ta}_3^2}{4G}\dd x+\F[\phi]\notag\\
 &+\int_\Gamma\left[
       \frac{Z^2}{2GI}+\frac G2\int\frac{\rho^2}{h}\dd n
      \right]\dd s
   +\int_{\Omega\setminus A_\Gamma}\frac{Ga^2}{2h}\dd x .
 \label{eq:hybridexact}
\end{align}
The dissipation is partitioned analogously using \eqref{eq:rhodiss}
and \eqref{eq:frictionexcess}. This is an exact representation change under
the stated geometric assumptions. There is no overlapping extra cohesive
energy and no second frictional mechanism.

Newly reconstructed surface records inherit $s_e=\int a\dd n$ and the
corresponding rate/history moments. They are not assigned zero histories
merely because their indices are new.
If an intact profile has $I=0$, it has no normalized surface row:
activation is handled by the bulk mixed law and its one-sided limit.
When $I>0$, the surface variables can be extracted from the existing bulk
state. Phase irreversibility remains attached to the bulk material phase,
not to the lifetime of a surface vertex.

The accumulated signed slip and accumulated absolute slip must be
distinguished: their rates are $\int w\dd n$ and $\int|w|\dd n$.
They agree in magnitude only under common-sign local motion.
No RSF-state transfer law follows from this derivation.

\section{What can be carried into the current discretization}

\subsection{One coherent P4 choice}

For the first fixed-direction growing-tip experiment, select the
\emph{bulk-authoritative} law \eqref{eq:bulkenergy}--\eqref{eq:phaseVI}.
The orientation $\Sg$ is prescribed in intact material as well as damaged
material. It is a candidate shear direction, not a pre-existing slipping
surface. The bulk equations then apply ahead of any reconstructed endpoint.

Retain independent DG $\T$ and a bulk representation of $a$ (or an equivalent
mixed local update); retain whatever physical slip history is needed for the
chosen frictional law. The surface can initially be an observer.
If surface reduction is later used mechanically, retain $\rho$ or explicitly
accept and measure the approximation of setting it to zero.
This is a governing-model change for the cohesive propagation route, not
an automatic modification of the frozen mature BP3 branch.

Alternatively, one may keep the current collective surface law as the
primary physical model and derive a crack-front/configurational force for
surface advance. That is a distinct choice. Appending surface records
inside old damage, as in P3, does not derive that force.

\subsection{Consistent discrete normalization}

The continuum identity $\int\chi\dd n=1$ must not be silently substituted
for a projected discrete identity. If the implemented nodal normalization is
\begin{equation}
 I_h=P\,i(\phi_h),\qquad P=M_\Gamma^{-1}B_\Gamma,
\end{equation}
its first and second derivatives and transpose actions are part of the
discrete energy. In general a projected $I_h(s)$ differs from the raw
column integral $i(s)$, so
$\int h(\phi_h)/I_h\dd n=i(s)/I_h(s)$ need not equal one pointwise.
The continuous profile energy and moment formulas are therefore not exact
proofs of an arbitrary projected scheme.

An exact finite-dimensional cohesive decomposition is available.
At one profile, let $W$ be positive quadrature weights, $H=\operatorname{diag}(h_i)$,
$\bm a$ its cohesive shear samples and $\bm1$ the vector of ones.
On positive-$h$ entries put
\begin{equation}
 I_d=\bm1^TWH\bm1,\qquad
 s_d=\bm1^TW\bm a,\qquad
 \bm\chi_d=H\bm1/I_d.
\end{equation}
Then, with $\bm\rho=\bm a-\bm\chi_d s_d$,
\begin{equation}
 \frac G2\bm a^TWH^{-1}\bm a
 =\frac{G s_d^2}{2I_d}
    +\frac G2\bm\rho^TWH^{-1}\bm\rho,\qquad
 \bm1^TW\bm\rho=0.
 \label{eq:discretevariance}
\end{equation}
At intact entries enforce $a_i=0$ or use the mixed form.
For multiple coupled surface degrees of freedom, derive the constrained
minimum with the actual interpolation and mass matrices.
Changing normalization merely to enforce \eqref{eq:discretevariance} would
change the existing discrete model; it requires a deliberate derivation,
not an unexplained numerical repair.

Whichever discrete energy is chosen, mechanics, phase residual, tangent and
work diagnostics must differentiate that same energy. A local driving
source cannot simply be added to the old surface phase row while retaining
unmodified mechanics and storage.

\section{Limits of the flat-profile identities}

\subsection{Curvature and normal heterogeneity}

For a curved tube $\bm X(s,n)$ with Jacobian $J(s,n)>0$,
$\dd x=J\dd s\,\dd n$.
If physical slip is defined by the unweighted normal integral
$s_e=\int a\dd n$ and the old $\chi=h/I$ is retained, then
\begin{equation}
 E_{c,\mathrm{profile}}
 =\frac{G s_e^2}{2I^2}\int J h\dd n ,
\end{equation}
which is not generally $G s_e^2/(2I)$.
For a symmetric profile and the usual planar-curve Jacobian
$J=1-\kappa n$, the odd contribution can cancel, but that is an additional
symmetry assumption and can fail at boundaries.

More generally, minimizing $\int JG(n)a^2/(2h)\dd n$ subject to
$\int a\dd n=s_e$ yields
\begin{equation}
 C_\Gamma=\int\frac{h}{J G(n)}\dd n,\qquad
 a=\frac{h}{J G(n)}\frac{s_e}{C_\Gamma},\qquad
 E_c=\frac{s_e^2}{2C_\Gamma}.
 \label{eq:curved}
\end{equation}
The multiplier is conjugate to the specified physical slip, and $Jq$,
not necessarily $q$, is constant in this constrained minimum.
For $J=1$ and constant $G$, $C_\Gamma=I/G$ and the earlier result follows.
One must also transform the weak traction, kinematics and dissipative
operators with the same measure. This note does not claim a general curved,
moving-surface reduction from the flat formula.

\subsection{Moving geometry, saturation, and RSF}

If geometry or orientation moves, differentiating the extension produces
transport, orientation and metric terms. The fixed-geometry source
$-s_e\dot\chi$ in \eqref{eq:rhoevolution} is only one part of that product rule.
Rigid rotations require a consistent objective tensor rate and rotated
kinematics; none is implicitly included here.

At $\phi\uparrow1$, $h$ diverges for \eqref{eq:h}; a fully failed limit must
be derived rather than capped. RSF introduces rate and state dependence and
an additional history law. Positivity of $\mu p_N|V|$ alone does not derive
the state-dependent energy, activation, or radiation terms of such an
extension. These issues remain outside the present constant-friction P4
derivation.

\section{Conclusions and analytical verification}

The derivation establishes the following concrete results:
\begin{enumerate}
 \item A bulk-defined free energy and dissipation potential yield explicit
       Maxwell, Coulomb and irreversible phase equations with the power
       balance \eqref{eq:balance}.
 \item Intact material has zero cohesive deformation but a finite limiting
       traction. The virgin activation condition follows from relaxed energy
       and friction, with no existing surface association or $I$ floor.
 \item The current surface storage is the normal-profile reduction of bulk
       cohesive storage. Its use requires explicit assumptions on
       representation, geometry and material coefficients.
 \item A surface moment is not generally a closed history when the normal
       profile evolves. Equations \eqref{eq:variance},
       \eqref{eq:rhodiss}, \eqref{eq:rhoevolution} and
       \eqref{eq:frictionexcess} quantify exactly what is retained or omitted.
 \item Keeping bulk fields authoritative and using
       \eqref{eq:hybridexact} connects the descriptions without duplicate
       energy or dissipation. An existing finite surface need not gate
       mechanical damage ahead of its endpoint.
\end{enumerate}

Four scalar identities were checked by symbolic algebra: the mechanical
power identity in transformed variables; the virgin minimized potential;
its compliance derivative; and the frictional stress interpolation.
One hundred deterministic random positive-weight profiles checked the
cohesive energy decomposition, cohesive viscous dissipation decomposition,
and the phase chain rule. Maximum absolute differences in nondimensional
test values were respectively $6.94\times10^{-18}$,
$2.78\times10^{-17}$ and $1.74\times10^{-18}$.
Centered finite differences of the reduced profile energy, with step
$10^{-5}$, differed from its analytic directional derivative by at most
$1.59\times10^{-11}$.
These checks verify algebra and differentiation, not a finite-element
propagation solution, mesh convergence, or constitutive calibration.

\appendix
\section{A compact derivation of the transformed power identity}

In scalar shear let $\tau=T-Ga$, $q=Ga/h$ and
$\psi=\tau^2/(2G)+Ga^2/(2h)$. Then
\begin{align}
 \dot\psi
 &=\frac{\tau}{G}\dot T+(q-\tau)\dot a
      -\frac{q^2}{2G}\dot h\\
 &=\frac{\tau}{G}(G\dot\gamma-T/\lambda)
   +(q-\tau)(w-a/\lambda)-\frac{q^2}{2G}\dot h\\
 &=\tau\dot\gamma-\frac{\tau^2}{\eta}
      -\frac{h q^2}{\eta}-(\tau-q)w
      -\frac{q^2}{2G}\dot h .
\end{align}
The scalar bulk dissipation is $\tau^2/\eta$ because the full tensor
contains two equal shear components. This confirms the factors of two in
\eqref{eq:balance} and avoids confusing engineering shear with
$\varepsilon_{xy}$.

\section{Practical comparison with the supplied reports}

The P4a report establishes zero mechanical variation for perturbations
wholly outside the represented strip, while its full normalization pullback
and independent high-order integration agree. Nothing here reinterprets
that result as an integration failure.

P3's zero-time map preserves physical stress approximately after a prescribed
geometry change and separately accounts for storage/load differences.
It does not supply virgin activation or justify resetting histories on a
newly reconstructed segment. The exact residual energy in
\eqref{eq:variance} gives an additional diagnostic for any proposed
bulk-to-surface reduction.

P2's evolving-profile energy and phase derivative are recovered on their
selected profile manifold, but equivalence to arbitrary bulk local evolution
is not established by that recovery. Existing P2 accuracy and branch
limitations, and all protected mature BP3 results, remain separate evidence.

\begin{thebibliography}{9}
\bibitem{fei2020}
F. Fei and J. Choo,
\emph{A phase-field model of frictional shear fracture in geologic materials},
Computer Methods in Applied Mechanics and Engineering 369 (2020), 113265.
\href{https://doi.org/10.1016/j.cma.2020.113265}{doi:10.1016/j.cma.2020.113265}.
Author manuscript:
\url{https://arxiv.org/abs/2003.04779}.

\bibitem{fei2026}
F. Fei, M. S. Mia, A. E. Elbanna and J. Choo,
\emph{A Phase-Field Model for Quasi-Dynamic Rupture Nucleation and
Propagation of In-Plane Faults},
International Journal for Numerical and Analytical Methods in Geomechanics
(2026).
\href{https://doi.org/10.1002/nag.70416}{doi:10.1002/nag.70416}.

\bibitem{project}
\raggedright
Project inputs reviewed for this derivation:
\texttt{P4A\_TIP\_REPORT.md} (19 October 2026);
\texttt{P3\_NATIVE\_COMPLETION\_REPORT.md};
\texttt{particle\_free\_history.tex}, particularly the evolving-profile
product rule and cohesive energy;
\texttt{PF6A\_FORMULATION.md}; and
\texttt{P2\_coupled\_branch\_equations.md}.
The inspected repository baseline was commit
\texttt{612b75d3570d81763407875f3a12c12999b9fa8e};
the supplied later reports describe newer work.
This note changes no repository source or production assumption by itself.
\end{thebibliography}
\end{document}
